\chapter{System overview}
\label{cap:visao}

\section{Executive summary}

The Rust VGA Kernel is an x86\_64 ELF executable loaded by GRUB. GRUB uses
the Multiboot2 protocol to transfer control to a small assembly-language
\emph{stub}, linked at 1 MiB. This stub builds page tables, enables PAE and
long mode and jumps to the Rust code, linked in the higher half at
\code{0xFFFFFFFF80000000}.

The Rust code initializes, in this order, the VGA renderer, the COM1 serial
port, the TTY layer with eight virtual consoles, the frame allocator, the GDT
with ring 3/TSS, the PS/2 keyboard and the IDT with PIC/PIT. It then enables
interrupts, prints a \emph{banner}, validates the Multiboot2 signature and
enters the shell, which displays the prompt \code{rinux>} on the active
console and processes commands read from the keyboard or from the serial
port.

The entire system fits in about 95 KiB of ELF and 3,223 lines of Rust
(2,477 effective lines of code, without comments and blank lines). There is no
dynamic heap allocation, scheduler or filesystem; there is an IDT, timer,
\emph{demand paging}, eight virtual consoles, four system calls and a
\code{getty}/login running in ring 3.

\section{Features}

\begin{table}[H]
\centering
\caption{Implemented features and their status.}
\label{tab:funcionalidades}
\begin{tabularx}{\linewidth}{@{}lXc@{}}
\toprule
\textbf{Area} & \textbf{Feature} & \textbf{Status} \\
\midrule
Boot & Valid Multiboot2 header & Complete \\
Boot & 32 $\rightarrow$ 64 bit transition & Complete \\
Memory & Static PML4/PDPT/PD, 1 GiB & Complete \\
Memory & Kernel higher half & Complete \\
Video & VGA text console 80$\times$25 & Complete \\
Video & Scrolling, cursor, tab, backspace & Complete \\
Video & Eight virtual consoles (\code{tty0}..\code{tty7}) & Complete \\
Video & Alt+F1..Alt+F8 and \code{tty} switching & Complete \\
Video & Graphics framebuffer / GOP & Absent \\
Keyboard & PS/2 8042, set 1, polling & Complete (subset) \\
Keyboard & Shift, Alt, Enter, Tab, Backspace & Complete \\
Keyboard & Alt+F1..Alt+F8 console selection & Complete \\
Keyboard & IRQ1, Caps Lock, AltGr, arrow keys & Absent \\
Serial & COM1 16550, 38400 8N1 & Complete (subset) \\
Serial & Input through the QEMU terminal & Complete \\
Shell & \code{help}, \code{clear}, \code{echo} & Complete \\
Shell & \code{info}, \code{tty}, \code{getty} & Complete \\
Shell & \code{reboot}, \code{halt} & Complete \\
System & Panic handler & Complete \\
System & IDT, exceptions and IRQ0 (timer) & Complete (subset) \\
System & GDT/TSS, ring 3 and U/S protection & Complete (demo) \\
System & Syscalls \code{exit/write/read/get\_tty} & Complete (subset) \\
User & \code{getty}/login in ring 3 & Complete (demo) \\
Memory & Frames, 4 KiB pages and demand paging & Complete (subset) \\
System & Heap, processes and scheduler & Absent \\
\bottomrule
\end{tabularx}
\end{table}

\section{Requirements}

\subsection{Hardware}

\begin{itemize}
  \item x86\_64-compatible PC with legacy BIOS or CSM enabled.
  \item PS/2 8042 keyboard controller (or equivalent emulation).
  \item VGA adapter in text mode, with a buffer at \code{0xB8000}.
  \item CD/DVD drive or equivalent device to boot the ISO.
  \item Enough memory to load the kernel and the ISO; 128 MiB is sufficient
        in QEMU.
\end{itemize}

\subsection{Development software}

\begin{itemize}
  \item Stable Rust with the \code{x86\_64-unknown-none} target installed.
  \item \code{rustup} (recommended; the Makefile prioritizes its toolchain).
  \item \code{ld.lld} (LLVM linker).
  \item \code{grub-mkrescue}, \code{xorriso}, \code{libisoburn} and
        \code{mtools}.
  \item QEMU (\code{qemu-system-x86\_64}) for execution and tests.
  \item For this documentation: a TeX Live distribution with
        \code{pdflatex} and \code{latexmk}.
\end{itemize}

\section{Directory structure}

\begin{lstlisting}[style=shell]
.
|-- Cargo.toml              metadados e perfis de compilacao
|-- Cargo.lock              lockfile (sem dependencias externas)
|-- rust-toolchain.toml     canal stable + alvo x86_64-unknown-none
|-- .cargo/config.toml      alvo padrao, linker e flags
|-- linker.ld               layout do ELF (stub baixo + kernel alto)
|-- Makefile                build do kernel, ISO e execucao no QEMU
|-- boot/grub/grub.cfg      menu Multiboot2 e console VGA
|-- src/
|   |-- main.rs             entrada Rust, macros e panic handler
|   |-- boot.rs             Multiboot2, GDT, paginacao e modo longo
|   |-- arch.rs             portas de I/O, interrupcoes, halt e reboot
|   |-- gdt.rs              GDT definitiva, TSS e pilhas de privilegio
|   |-- idt.rs              IDT, stubs, PIC, PIT e porta int 0x80
|   |-- paging.rs           frames, tabelas 4 KiB e demand paging
|   |-- user.rs             ring 3, demos e getty/login
|   |-- syscall.rs          exit/write/read/get_tty
|   |-- vga.rs              renderizador VGA 80x25
|   |-- tty.rs              consoles virtuais tty0..tty7
|   |-- keyboard.rs         teclado PS/2 (8042)
|   |-- serial.rs           UART 16550 / COM1
|   |-- shell.rs            prompt e comandos
|   `-- sync.rs             spinlock do console
|-- docs/                   esta documentacao (LaTeX e PDF)
|-- build/                  arvore temporaria da ISO (gerada)
|-- target/                 artefatos do Cargo (gerado)
|-- vga-kernel.iso          imagem inicializavel (gerada)
`-- README.md               guia rapido de uso
\end{lstlisting}

\section{Code metrics}

\begin{table}[H]
\centering
\caption{Metrics of the Rust modules (count by physical lines).}
\label{tab:metricas}
\begin{tabular}{@{}lrrl@{}}
\toprule
\textbf{Module} & \textbf{Lines} & \textbf{Functions} & \textbf{Responsibility} \\
\midrule
\code{boot.rs}     & 202 & 0 & Boot stub and page tables \\
\code{keyboard.rs} & 271 & 10 & PS/2 8042 driver \\
\code{vga.rs}      &  74 & 5 & VGA text renderer \\
\code{tty.rs}      & 448 & 31 & Virtual consoles and queues \\
\code{shell.rs}    & 245 & 9 & Command line and demos \\
\code{main.rs}     & 119 & 2 & Rust entry point, macros, panic \\
\code{serial.rs}   &  70 & 5 & UART 16550 / COM1 \\
\code{syscall.rs}  & 129 & 3 & exit/write/read/get\_tty \\
\code{arch.rs}     &  69 & 5 & I/O ports, \code{sti}, halt, reboot \\
\code{sync.rs}     &  56 & 5 & Spinlock \\
\code{gdt.rs}      & 159 & 7 & GDT, TSS and privilege stacks \\
\code{idt.rs}      & 504 & 14 & IDT, PIC, PIT, exceptions and IRQs \\
\code{paging.rs}   & 559 & 23 & Frames, tables and demand paging \\
\code{user.rs}     & 318 & 10 & Ring 3 entry, demos and getty \\
\midrule
\textbf{Total}     & \textbf{3,223} & \textbf{129} & 2,477 effective lines of code \\
\bottomrule
\end{tabular}
\end{table}

\begin{nota}[On the counting]
The count includes 351 comment lines and 395 blank lines. The number
of functions includes \code{impl} methods, \code{const} functions
and trait functions (\code{Deref}, \code{DerefMut}, \code{Drop} and
\code{fmt::Write}); it does not include \code{extern} declarations without a
body. \code{boot.rs} has no Rust functions because its logic is in
\code{global\_asm!}; the \emph{stub} is counted as an assembly block.
\end{nota}

\section{Summary of the execution flow}

\begin{enumerate}
  \item The BIOS selects the CD and executes El Torito.
  \item GRUB loads the ELF and reads the Multiboot2 header.
  \item GRUB enters 32-bit protected mode and jumps to \code{\_start}.
  \item The stub saves \code{eax} and \code{ebx}, builds the pages, loads the
        GDT and enables long mode.
  \item The stub jumps to \code{kernel\_main(magic, info)} in the higher half.
  \item Rust clears the VGA screen, configures the serial port, initializes
        the TTY layer with eight virtual consoles, initializes the frame
        allocator and the GDT/TSS, initializes the keyboard and installs the
        IDT with PIC/PIT.
  \item Rust enables interrupts and the shell displays \code{rinux>} on the
        active console, processing commands until \code{halt}, \code{reboot}
        or a physical reset.
\end{enumerate}
